\documentclass[11pt,letterpaper]{article}
\usepackage[margin=1in]{geometry}
\usepackage[T1]{fontenc}
\usepackage{amsmath,amssymb,amsthm,booktabs,array,graphicx,float}
\usepackage[nopatch={eqnum,footnote}]{microtype}
\usepackage{setspace}
\usepackage[numbers,square,sort&compress]{natbib}
\usepackage[hidelinks]{hyperref}
\usepackage{lineno}
\usepackage[font=small,labelfont=bf,labelsep=period,justification=raggedright,singlelinecheck=off]{caption}
\usepackage{fancyhdr}
\pagestyle{fancy}
\fancyhf{}
\fancyfoot[C]{\thepage}
\renewcommand{\headrulewidth}{0pt}
\newcommand{\E}{\mathbb E}
\newcommand{\R}{\mathbb R}
\newcommand{\norm}[1]{\left\lVert#1\right\rVert}
\newtheorem{proposition}{Example}
\setlength{\emergencystretch}{2em}
\setlength{\parindent}{0pt}
\setlength{\parskip}{0.4em}
\widowpenalty=10000
\clubpenalty=10000
\raggedright
\doublespacing
\hypersetup{pdftitle={Supplementary Information: Repeatable interaction scores do not guarantee a correct circuit},pdfauthor={Ihor Kendiukhov}}
\begin{document}
\linenumbers
{\Large\textbf{Supplementary Information}}

Repeatable interaction scores do not guarantee a correct circuit

Ihor Kendiukhov

\renewcommand{\thetable}{S\arabic{table}}
\section*{Which claims passed their planned tests?}
\begin{table}[H]
\singlespacing\centering\small
\caption{The main questions fixed before each test and what the results support.
A positive part of a test does not turn a failed full test into a pass.}
\begin{tabular}{>{\raggedright\arraybackslash}p{0.16\linewidth}>{\raggedright\arraybackslash}p{0.36\linewidth}>{\raggedright\arraybackslash}p{0.37\linewidth}}
\toprule
Study & Question fixed in advance & What passed or failed\\
\midrule
E31.1 & Do a few corrected block-pair scores repeat? & Main test failed;
uncorrected scores suggested later tests.\\
E31.2 & Do a few head pairs repeat and give a better choice of parts? & Full
test failed; score coverage and reproduction of block interactions passed.\\
E31.2 later check & Does a product of part scores explain the table? & Designed
after the main result; no new model runs.\\
E32.1 & Do pair scores find direct links better than simpler methods? & Main
test failed; longer-path and product-fit checks were added afterward.\\
E32.2 & Does choosing parts first help find direct links? & Both models failed
the full requirements; all results retained.\\
E32.3 & Do small part groups work, and is the selector better? & Groups passed
3/3; the selection method beat comparisons in 1/3.\\
E32.4 & Does a small connection set preserve the computation? & Full test
passed in 1/3 models.\\
E33 & Do scores repeat, concentrate in few pairs, and emphasize longer paths? &
Repeatability passed 3/3; concentration and product fit passed 2/3;
indirect pairs scored higher in both models with longer paths.\\
\bottomrule
\end{tabular}
\end{table}

\section*{Tests aimed at individual connections}
Cases 2 and 21 first chose a short list of parts, then targeted changes at
individual connections between them. We compared this with pair scores,
products of fitted part scores, and EAP methods tested across all connections.
Neither model met the full success requirements. Table~\ref{tab:directed}
reports both cases. We favor the comparison method by reporting whichever
specified EAP version scored best after evaluation. This does not tell us
which version to choose on a new task. The final column counts separate
query/key/value connections; AP instead scores links grouped by their whole
head or MLP. These columns therefore use different counting units.

\begin{table}[H]
\singlespacing\centering\footnotesize
\caption{Ranking direct links on new inputs after first choosing parts.
``Best EAP'' is the best-scoring specified EAP version for that model.
Neither case passes the full planned test.}
\label{tab:directed}
\begin{tabular}{lrrrrr}
\toprule
Case & Pair AP & Product AP & Link-change AP & Best EAP AP & Links tested \\
\midrule
2 & 0.717 & 0.717 & 0.189 & 0.339 & 18.0\% \\
21 & 0.700 & 0.417 & 0.591 & 0.417 & 16.0\% \\
\bottomrule
\end{tabular}
% AP is on promoted whole-component labels. Candidate fraction uses fine Q/K/V edges.
% Best EAP* variant is eap_absolute_max in case 2 and eap_ig_signed_max in case 21.
% Both cases fail the registered useful and superiority conjunctions.

\end{table}

\section*{Saved data and checks}
\subsection*{Files and score checks.}
The accompanying evidence file lists the original source files, their hashes,
the exact scores, and whether an analysis was planned or added after the
results. The numbers are read from saved experimental outputs rather than
rounded research notes. Original plans, model outputs, and score files are
unchanged. Figures and the manuscript are built from separate derived files.
The fresh-study registration is
\path{notes/E33-SPARSE-INTERACTION-DIAGNOSTICS-REGISTRATION.json};
its complete score record is
\path{runs/e33_sparse_interaction_diagnostics_v1/summary.json}.
An independent internal check recomputed all six new pair tables and scores
exactly from saved outputs. It also checked input identities, reference links,
source-code and model hashes, and run counts. This checks the saved executions;
it is not an external replication or a second collection of model responses.

\subsection*{Changes to the execution plans.}
The first block-pair study stopped because two file mappings used a different
key order. A separately recorded repair reused discovery data and was fixed
before measuring confirmation changes. The small-part-group study also used a
documented input-repair version. These histories are saved with the experiments.
Reusing data already seen is not counted as a new independent confirmation.
\end{document}
